\documentclass[border=5mm]{standalone}
\usepackage[UTF8,fontset=fandol]{ctex}
\usepackage{amsmath}
\usepackage{pgfplots}
\pgfplotsset{compat=1.18}
\definecolor{firstnetwork}{HTML}{83C9C4}
\definecolor{secondnetwork}{HTML}{C67D60}
\definecolor{composednetwork}{HTML}{526A77}
% Independent schematic solution to UDL Problem 4.1.
% Source: Fig.4.8, printed54/PDF68; checked against answer Fig.4.2, printed14/PDF14.
% The book does not give all exact knot coordinates.
% All rational coordinates below are a declared approximate reconstruction.
% f vertices: (-1,0),(-2/5,1),(0,-1),(3/5,1),(1,0).
% g vertices: (-1,3/5),(-1/5,-4/5),(1/2,-2/3),(4/5,0),(1,0).
% The composite uses exact inverses of the knots on these particular affine pieces.
% 15 candidate vertices reduce to 13 after two collinear plateau vertices are removed.
% This is 12 maximal affine pieces for this representative, not a theorem about
% arbitrary three-unit networks or a claim of recovered original parameters.
\begin{document}
\begin{tikzpicture}
\pgfplotsset{
  schematic/.style={
    xmin=-1,xmax=1,ymin=-1,ymax=1,
    xtick={-1,0,1},ytick={-1,0,1},
    tick label style={font=\small},
    label style={font=\small},
    title style={font=\small},
    axis line style={black,thin},tick style={black,thin},
    axis on top,clip=true,
    every axis plot/.append style={line width=1.2pt},
  }
}
\begin{axis}[schematic,name=first,width=66mm,height=54mm,
  title={(a) $y=f(x)$},xlabel={$x$},ylabel={$y$}]
\addplot[gray!50,densely dotted,thin] coordinates {(-1,0)(1,0)};
\addplot[gray!50,densely dotted,thin] coordinates {(0,-1)(0,1)};
\addplot[firstnetwork] coordinates {(-1,0)({-2/5},1)(0,-1)({3/5},1)(1,0)};
\end{axis}
\begin{axis}[schematic,name=second,at={(78mm,0)},anchor=south west,
  width=66mm,height=54mm,title={(b) $y'=g(y)$},xlabel={$y$},ylabel={$y'$}]
\addplot[gray!50,densely dotted,thin] coordinates {(-1,0)(1,0)};
\addplot[gray!50,densely dotted,thin] coordinates {(0,-1)(0,1)};
\addplot[secondnetwork] coordinates
  {(-1,{3/5})({-1/5},{-4/5})({1/2},{-2/3})({4/5},0)(1,0)};
\end{axis}
\begin{axis}[schematic,name=composite,at={(0,-72mm)},anchor=south west,
  width=144mm,height=63mm,title={(c) $y'=g(f(x))$},xlabel={$x$},ylabel={$y'$}]
\addplot[gray!50,densely dotted,thin] coordinates {(-1,0)(1,0)};
\addplot[gray!50,densely dotted,thin] coordinates {(0,-1)(0,1)};
\addplot[composednetwork] coordinates {
  (-1,{-16/21})
  ({-7/10},{-2/3})
  ({-13/25},0)
  ({-9/25},0)
  ({-3/10},{-2/3})
  ({-4/25},{-4/5})
  (0,{3/5})
  ({6/25},{-4/5})
  ({9/20},{-2/3})
  ({27/50},0)
  ({17/25},0)
  ({4/5},{-2/3})
  (1,{-16/21})
};
\end{axis}
\node[anchor=north,align=center,text width=145mm,font=\small]
  at ([yshift=-13mm]composite.south) {
  习题 DL019：两个折线网络的复合。\\
  本图节点按原图估读后选为便于复算的有理数；不表示原书给定了这些精确参数。
};
\end{tikzpicture}
\end{document}
